% Part: sets-functions-relations
% Chapter: functions
% Section: partial-functions

\documentclass[../../../include/open-logic-section]{subfiles}


\begin{document}

\olfileid[ar]{sfr}{fun}{par}

\olsection{الدوال الجزئية}

\begin{explain}
من المفيد أحيانًا تخفيف شروط تعريف الدالة بحيث لا يُشترط أن يكون
مخرجها معرّفًا لكل مدخل ممكن. وتسمى التعيينات من هذا القبيل
\emph{دوالًا جزئية}.
\end{explain}

\begin{defn}
إن \emph{الدالة الجزئية} $f \colon A \pto B$ تعيينٌ يرسل كل
!!{element} من~$A$ إلى !!{element} واحد على الأكثر من~$B$.
إذا عيّنت $f$ عنصرًا من~$B$ للمدخل $x \in A$، نقول إن قيمة $f(x)$
\emph{معرّفة}، وإلا فهي \emph{غير معرّفة}. وإذا كانت قيمة $f(x)$ معرّفة،
نكتب $f(x) \fdefined$، وإلا نكتب $f(x) \fundefined$. أما
\emph{مجال} الدالة الجزئية~$f$ فهو المجموعة الجزئية من~$A$ المؤلفة
من جميع العناصر التي تكون الدالة معرّفة عندها؛ أي $\dom{f} = \Setabs{x \in A}{f(x) \fdefined}$.
\end{defn}

\begin{ex}
كل دالة $f\colon A \to B$ هي دالة جزئية أيضًا. أما الدوال
الجزئية المعرّفة عند كل عنصر من~$A$---أي ما سميناه حتى الآن
دالة فحسب---فتسمى أيضًا دوالًا \emph{كلية}.
\end{ex}

\begin{ex}
الدالة الجزئية $f \colon \Real \pto \Real$ المعطاة بـ$f(x) = 1/x$
غير معرّفة عند $x = 0$، ومعرّفة في كل موضع آخر.
\end{ex}

\begin{prob}
إذا أُعطيت الدالة الجزئية $f\colon A \pto B$، فعرّف الدالة الجزئية $g\colon B \pto
A$ كما يأتي: لكل $y \in B$، إذا وُجد عنصر وحيد $x \in A$ بحيث
$f(x) = y$، فليكن $g(y) = x$؛ وإلا فـ$g(y) \fundefined$. بيّن أنه
إذا كانت $f$ متباينة، فإن $g(f(x)) = x$ لكل $x \in \dom{f}$، وكذلك
$f(g(y)) = y$ لكل $y \in \ran{f}$.
\end{prob}

\begin{defn}[الرسم البياني لدالة جزئية]
لتكن $f\colon A \pto B$ دالة جزئية. إن \emph{الرسم البياني} لـ~$f$
هو العلاقة $R_f \subseteq A \times B$ المعرّفة بـ
\[
R_f = \Setabs{\tuple{x,y}}{f(x) = y}.
\]
\end{defn}

\begin{prop}
افترض أن $R \subseteq A \times B$ تتصف بالخاصية الآتية: متى صدق $Rxy$
و$Rxy'$ كان $y = y'$. عندئذ تكون $R$ الرسم البياني للدالة الجزئية
$f\colon A \pto B$ المعرّفة كما يأتي: إذا وُجد $y$ بحيث
$Rxy$، كان $f(x) = y$؛ وإلا فـ$f(x) \fundefined$. وإذا كانت $R$
\emph{تسلسلية} أيضًا، أي إذا كان لكل $x \in A$ عنصر $y \in B$ بحيث
$Rxy$، فإن $f$ كلية.
\end{prop}

\begin{proof}
افترض أن ثمة $y$ بحيث $Rxy$. ولو وُجد $y'
\neq y$ آخر بحيث $Rxy'$، لانتهكت الخاصية المفروضة على~$R$.
ومن ثم، إذا وُجد $y$ بحيث $Rxy$، كان ذلك $y$
وحيدًا، فتكون $f$ معرّفة تعريفًا سليمًا. ومن الواضح أن $R_f = R$ وأن $f$
كلية إذا كانت~$R$ تسلسلية.
\end{proof}

\end{document}
